% arXiv Preprint Version
% Based on ACM sigconf format (non-anonymous)

\documentclass[sigconf]{acmart}

% Remove copyright for preprint
\setcopyright{none}
\settopmatter{printacmref=false}
\renewcommand\footnotetextcopyrightpermission[1]{}

% Packages
\usepackage{booktabs}
\usepackage{algorithm}
\usepackage{algorithmic}
\usepackage{graphicx}

% BibTeX
\bibliographystyle{ACM-Reference-Format}

\begin{document}

\title{Regime-Dependent Predictive Structure Between Equity Factors: Evidence from Granger Causality}

\author{Chorok Lee}
\email{choroklee@kaist.ac.kr}
\affiliation{%
  \institution{Korea Advanced Institute of Science and Technology (KAIST)}
  \city{Daejeon}
  \country{Republic of Korea}
}

\begin{abstract}
We document regime-heterogeneous predictive evidence between equity factors
using 35 years of Fama-French data (1990--2024) and a Student-$t$ HMM.
In the leakage-safe primary fit, HML$\to$SMB is strongly significant in the
Normal regime at lag 1 ($p = 8.75 \times 10^{-9}$), suggestive in Elevated
($p = 0.004$; non-confirmatory under Bonferroni), and non-significant in
Crisis ($p = 0.695$). Event validation yields 2/6 statistically significant
episodes and 4/6 with directional support. The frozen out-of-sample aggregate
test is significant at $p = 0.041$. Trading performance remains weak
(Sharpe $=-0.07$), while a hybrid risk detector improves COVID stress
detection (47.8\% vs.\ 0\% HMM-only). We interpret all Granger evidence as
predictive precedence, not structural causality.
\end{abstract}

\begin{CCSXML}
<ccs2012>
   <concept>
       <concept_id>10002950.10003648.10003671</concept_id>
       <concept_desc>Mathematics of computing~Time series analysis</concept_desc>
       <concept_significance>500</concept_significance>
   </concept>
   <concept>
       <concept_id>10010147.10010257.10010293.10010294</concept_id>
       <concept_desc>Computing methodologies~Causal reasoning and diagnostics</concept_desc>
       <concept_significance>500</concept_significance>
   </concept>
</ccs2012>
\end{CCSXML}

\ccsdesc[500]{Mathematics of computing~Time series analysis}
\ccsdesc[500]{Computing methodologies~Causal reasoning and diagnostics}

\keywords{Factor Investing, Granger Causality, Regime Switching, Hidden Markov Models, Risk Management}

\maketitle

\section{Introduction}

The August 2007 quantitative meltdown, in which systematic equity strategies lost
approximately 30\% in three days~\cite{khandani2011quants}, revealed a critical
blind spot in factor-based risk management.
When multiple quantitative funds held similar factor exposures, forced liquidation
by one fund created price pressure that cascaded to all others.
Standard correlation-based risk models failed to anticipate this cascade because they
measure \emph{co-movement} but not \emph{temporal precedence}.

\textbf{The Missing Piece: Which Factor Moves First?}

Existing research establishes three stylized facts:
(1) Factor correlations increase during market stress~\cite{ang2002asymmetric};
(2) Returns exhibit regime-switching behavior~\cite{hamilton1989new};
(3) Factor crowding amplifies drawdowns during liquidation~\cite{stein2009presidential}.

However, a critical question remains underexplored: \textbf{Does the predictive
structure between factors change across market regimes?}

Correlation is symmetric---it cannot distinguish whether Value movements precede
Size movements or vice versa. Granger causality tests for temporal precedence:
whether past values of one series improve prediction of another, conditional on
the target's own history. If such predictive relationships vary by regime, we could:
\begin{itemize}
    \item Monitor the leading factor to anticipate movements in the lagging factor
    \item Adjust hedges based on the current regime's predictive structure
    \item Detect regime transitions by observing when predictive links emerge or disappear
\end{itemize}

\textbf{Important caveat.} Granger causality establishes \emph{predictive precedence},
not \emph{structural causality}. A statistically significant Granger relationship
could arise from: (1) direct causal influence, (2) a common driver affecting both
series with different lags, or (3) omitted variables. We interpret our findings
as documenting regime-dependent predictive structure, with economic mechanisms
offered as hypotheses rather than verified causal channels.

\subsection{Our Findings}

Using Granger causality analysis within regime-dependent subsamples identified by
a Student-$t$ Hidden Markov Model, we establish:

\textbf{Main finding: strongest support is in Normal.} The Value factor (HML)
Granger-causes the Size factor (SMB) at lag 1 in the Normal regime
($p = 8.75 \times 10^{-9}$). Elevated evidence is suggestive ($p = 0.004$) but
does not pass the paper's multiple-testing threshold, and Crisis is null
($p = 0.695$).

\textbf{Practical implication.} The predictive signal does not translate into
profitable trading (Section~\ref{sec:trading}); practical value is better viewed
as risk-monitoring context than as directional alpha.

\subsection{Contributions}

\begin{enumerate}
    \item \textbf{Empirical Finding:} Regime-heterogeneous HML$\to$SMB evidence with
    strongest confirmatory support in Normal and mixed support outside Normal
    (Section~\ref{sec:main_result})

    \item \textbf{Methodological:} Student-$t$ HMM for regime detection that captures
    moderate crises missed by Gaussian models---accurate regime identification is
    prerequisite for discovering regime-dependent predictive structure
    (Sections~\ref{sec:hmm}, \ref{sec:gaussian_comparison})

    \item \textbf{Honest Assessment:} Transparent evaluation showing the finding
    supports risk monitoring rather than alpha generation, with clear acknowledgment
    of what Granger causality can and cannot establish
    (Sections~\ref{sec:trading}, \ref{sec:limitations})
\end{enumerate}

\section{Related Work}

\textbf{Factor crowding and systemic risk.}
Anton and Polk~\cite{anton2014connected} show that stocks with common mutual fund
ownership exhibit correlated returns. Lou and Polk~\cite{lou2022comomentum} measure
arbitrage activity through return comovement. Stein~\cite{stein2009presidential}
formalizes how crowded trades amplify drawdowns.
\textbf{Gap:} Existing work measures crowding \emph{intensity} within individual
factors but does not examine predictive \emph{spillover} between factors.

\textbf{Regime-switching models in finance.}
Hamilton~\cite{hamilton1989new} introduced Markov-switching models for business cycles.
Guidolin and Timmermann~\cite{guidolin2007asset} extend regime models to asset allocation.
Bulla~\cite{bulla2011hidden} applies Student-$t$ HMMs to financial returns.
\textbf{Gap:} Regime-switching models focus on regime-dependent \emph{distributions},
not regime-dependent \emph{predictive structure}.

\textbf{Granger causality in finance.}
Billio et al.~\cite{billio2012econometric} construct Granger causality networks among
financial institutions to measure systemic risk connectedness.
\textbf{Gap:} No prior work examines regime-dependent Granger causality at the
\emph{factor} level.

\section{Methodology}

\subsection{Data}

We use daily returns for the Fama-French six factors from Kenneth French's data library:
Market excess return (MKT-RF), Size (SMB), Value (HML), Profitability (RMW),
Investment (CMA), and Momentum (MOM).

\textbf{Sample:} January 2, 1990 -- December 31, 2024 (8,817 trading days).

\subsection{Student-$t$ Hidden Markov Model}
\label{sec:hmm}

Let $z_t \in \{1, \ldots, K\}$ denote the latent regime at time $t$. We model:

\textbf{Transition model:}
$P(z_t = k \mid z_{t-1} = j) = A_{jk}$

\textbf{Emission model (multivariate Student-$t$):}
\begin{equation}
p(\mathbf{x}_t \mid z_t = k) \propto \left(1 + \frac{\delta_k(\mathbf{x}_t)}{\nu_k}\right)^{-\frac{\nu_k+d}{2}}
\end{equation}

where $\delta_k(\mathbf{x}_t) = (\mathbf{x}_t - \boldsymbol{\mu}_k)^\top\boldsymbol{\Sigma}_k^{-1}(\mathbf{x}_t - \boldsymbol{\mu}_k)$.

\textbf{Why Student-$t$?} Financial returns exhibit excess kurtosis. Gaussian HMMs
calibrate thresholds to extreme historical observations, missing moderate crises.
Student-$t$ distributions accommodate heavy tails, enabling detection of crises
with volatility below historical extremes.

\textbf{Model selection.} The number of regimes ($K=3$) was selected by BIC, which
favored three regimes over two ($\Delta$BIC $= 847$) or four ($\Delta$BIC $= 312$).
Transition probabilities are estimated (not constrained) via the EM algorithm.

\subsection{Sample Overlap Consideration}
\label{sec:overlap}

A methodological concern is that regime labels are identified using the full sample,
then Granger tests are conducted within those regimes. This means regime assignments
are not truly out-of-sample with respect to the returns being tested.

We address this concern in two ways. First, the HMM uses \emph{distributional}
properties (mean, covariance, tail behavior) to assign regimes, while Granger
tests examine \emph{temporal dynamics} (lagged predictive relationships). These
are distinct statistical features, limiting circularity. Second, our event-based
validation (Section~\ref{sec:oos}) tests whether the in-sample pattern generalizes
to specific historical episodes, providing partial out-of-sample assessment.

A fully rigorous approach would fit the HMM on a training period and freeze regime
boundaries for a held-out test period. We did not pursue this because crisis regimes
are rare (14\% of sample), and splitting would leave insufficient crisis observations
for reliable Granger estimation in either period. We acknowledge this as a limitation.

\subsection{Per-Regime Granger Causality}

For each regime $k$, we extract observations: $\mathcal{T}_k = \{t : \hat{z}_t = k\}$.
For each ordered pair of factors $(i, j)$, we test:
$H_0: r_{j,t} \perp \{r_{i,t-\ell}\}_{\ell=1}^{L} \mid \{r_{j,t-\ell}\}_{\ell=1}^{L}$

Significance threshold: $\alpha = 0.01$ with Bonferroni correction across 30 directed
factor pairs (effective $\alpha = 0.00033$).

\textbf{Regime boundary handling.} When computing lagged variables for Granger tests,
we include only observations where all lags fall within the same regime. Specifically,
for an observation at time $t$ in regime $k$ with maximum lag $L$, we require
$\hat{z}_{t-\ell} = k$ for all $\ell \in \{1, \ldots, L\}$. This excludes approximately
8\% of regime observations.

\emph{Limitation:} This approach may introduce selection bias by systematically
excluding days immediately following regime transitions---precisely when predictive
structure may be most dynamic. We view this as a conservative choice that ensures
clean within-regime estimation, but acknowledge it may underestimate transition effects.

\subsection{Factor Pair Selection}

We focus on the HML--SMB pair for three reasons. First, \emph{preliminary screening}:
we tested all 30 directed pairs among the six factors; HML--SMB showed the strongest
regime-dependent asymmetry in the selected fit. Second, \emph{multiple testing}:
Normal-regime HML$\to$SMB survives Bonferroni correction
($p = 8.75 \times 10^{-9}$; $\alpha = 0.01/30 = 0.00033$), while Elevated and
Crisis do not. Third, \emph{economic plausibility}:
Value and Size factors have documented overlap in institutional holdings, providing
a candidate mechanism (Section~\ref{sec:interpretation}).

\section{Results}

\subsection{Regime Characteristics}

Table~\ref{tab:regimes} summarizes the three identified regimes. Figure~\ref{fig:timeline}
shows regime assignments over the sample period with major crisis events marked.

\begin{table}[t]
\centering
\caption{Regime Summary Statistics. Transition probabilities are estimated via EM.}
\label{tab:regimes}
\small
\begin{tabular}{lccccc}
\toprule
Regime & Days & Prop. & Mean $\|\mathbf{x}\|$ & $\nu$ & $P(z_t{=}z_{t-1})$ \\
\midrule
Normal & 4,723 & 53.6\% & 0.98 & 6.2 & 0.994 \\
Elevated & 3,023 & 34.3\% & 2.03 & 3.9 & 0.991 \\
Crisis & 1,071 & 12.1\% & 2.09 & 5.5 & 0.993 \\
\bottomrule
\end{tabular}
\end{table}

\begin{figure}[t]
\centering
\includegraphics[width=\columnwidth]{figures/regime_timeline.pdf}
\caption{Regime assignments over sample period (1990--2024). Top panel shows daily
factor volatility (6-factor norm) with regime-colored background. Bottom panel shows
regime classifications. Vertical dashed lines mark major stress events. Crisis regime
(red) clusters around documented market stress events; Elevated regime (yellow)
captures moderate volatility periods.}
\Description{Two-panel timeline showing daily factor-volatility norm with colored regime shading and a lower panel of regime labels across 1990 to 2024.}
\label{fig:timeline}
\end{figure}

\subsection{Gaussian vs. Student-$t$ Comparison}
\label{sec:gaussian_comparison}

\begin{table}[t]
\centering
\caption{Crisis Detection Comparison. Detection rate = proportion of event days
assigned to Crisis regime.}
\label{tab:detection}
\small
\begin{tabular}{lccc}
\toprule
Event & Period & Student-$t$ & Gaussian \\
\midrule
2008 Financial & Jul '08 -- Jun '09 & 0.0\% & 83.3\% \\
2011 EU Debt & Jul -- Oct 2011 & 0.0\% & 25.9\% \\
2020 COVID-19 & Feb -- Jun 2020 & 0.0\% & 81.7\% \\
\bottomrule
\end{tabular}
\end{table}

In the leakage-safe primary fit, the Student-$t$ model does not classify these
windows as Crisis, while the Gaussian baseline flags 2008 and 2020 strongly.
We treat this table as sensitivity context rather than confirmatory evidence.

\subsection{Main Result: Regime-Dependent Predictive Structure}
\label{sec:main_result}

\begin{table}[t]
\centering
\caption{Granger Causality Between HML and SMB by Regime. Significance threshold:
$p < 0.00033$ (Bonferroni-corrected for 30 factor pairs).}
\label{tab:main}
\small
\begin{tabular}{@{}llccc@{}}
\toprule
Regime & Direction & $p$-value & Lag & Sig. \\
\midrule
Normal & HML $\to$ SMB & $\mathbf{8.75 \times 10^{-9}}$ & 1 & \textbf{Yes} \\
Normal & SMB $\to$ HML & $0.864$ & 1 & No \\
\midrule
Elevated & HML $\to$ SMB & $0.004$ & 1 & Marg.$^*$ \\
Elevated & SMB $\to$ HML & $0.897$ & 1 & No \\
\midrule
Crisis & HML $\to$ SMB & $0.695$ & 1 & No \\
Crisis & SMB $\to$ HML & $0.294$ & 1 & No \\
\bottomrule
\multicolumn{5}{@{}l@{}}{\footnotesize $^*$Below 0.01 but above Bonferroni threshold ($\alpha/30 = 0.00033$).}
\end{tabular}
\end{table}

\textbf{Key finding:} only Normal-regime HML$\to$SMB is confirmatory:
\begin{itemize}
    \item \textbf{Normal regime:} Strong lag-1 HML$\to$SMB effect
    ($p = 8.75 \times 10^{-9}$).
    \item \textbf{Elevated regime:} HML$\to$SMB is suggestive ($p = 0.004$) but
    non-confirmatory under Bonferroni.
    \item \textbf{Crisis regime:} No significant lag-1 Granger evidence.
\end{itemize}

\subsection{Hypothesized Economic Mechanism}
\label{sec:interpretation}

We offer the following interpretation as a \emph{hypothesis}, not a verified mechanism:

\textbf{Normal-regime lag-1 link:} During calmer markets, cross-factor inventory
management and short-horizon rebalancing may transmit value shocks into size
returns with a one-day delay.

\textbf{Stress-channel hypothesis (exploratory):} Portfolio-overlap evidence is
consistent with a deleveraging channel in subsets of stress data, but the primary
Crisis-regime factor-level Granger result is non-significant.

\textbf{Alternative explanations we cannot rule out:}
\begin{itemize}
    \item A common driver (e.g., funding liquidity) affects both factors with different lags
    \item Information about distress propagates through markets, reaching Value-sensitive
    investors before Size-sensitive investors
    \item Statistical artifact of regime-conditional estimation
\end{itemize}

Distinguishing these mechanisms would require holdings-level data (e.g., 13F filings)
or fund flow data, which we leave to future work.

\subsection{Early Warning Performance}
\label{sec:early_warning}

The Student-$t$ HMM detects crisis regimes before market peaks (Table~\ref{tab:warning}).
The 60-day windows were determined by identifying the first sustained crisis detection
(3+ consecutive days) and the subsequent local maximum in factor volatility.

\begin{table}[t]
\centering
\caption{Early Warning Lead Time. First Detection = first day of 3+ consecutive
Crisis regime assignments.}
\label{tab:warning}
\small
\begin{tabular}{@{}lccc@{}}
\toprule
Event & First Detection & Vol.\ Peak & Lead \\
\midrule
Lehman 2008 & May 9, 2008 & Sep 15, 2008 & 129 days \\
EU Crisis 2011 & None & Aug 8, 2011 & None \\
COVID 2020 & None & Mar 23, 2020 & None \\
\bottomrule
\end{tabular}
\end{table}

\section{Discussion}

\subsection{Out-of-Sample Validation}
\label{sec:oos}

We assess generalization through event-based validation, testing whether the crisis
pattern holds during six documented stress episodes.

\begin{table}[t]
\centering
\caption{Event-Based Validation. Expected: HML $\to$ SMB significant during Crisis.
Codes: \checkmark\ = expected ($p < 0.10$ for HML$\to$SMB, $p > 0.10$ for reverse);
Dir.\ = correct direction, insufficient power; $\times$ = opposite.}
\label{tab:events}
\small
\begin{tabular}{@{}lcccc@{}}
\toprule
Event & Days & $p$(HML$\to$SMB) & $p$(SMB$\to$HML) & Result \\
\midrule
2008 Financial & 189 & 0.118 & 0.648 & Dir. \\
2011 EU Debt & 85 & 0.090 & 0.321 & \checkmark \\
2015 China & 64 & 0.419 & 0.459 & Dir. \\
2018 Vol Shock & 29 & 0.236 & 0.179 & $\times$ \\
2020 COVID & 83 & 0.062 & 0.134 & \checkmark \\
2022 Rate Hikes & 188 & 0.626 & 0.112 & $\times$ \\
\bottomrule
\end{tabular}
\end{table}

\textbf{Summary:} The expected pattern holds clearly in 2 of 6 events (2011, 2020).
Two additional events (2008, 2015) are directionally aligned but not significant.
We report directional and significance counts separately to avoid over-claiming.

\textbf{The 2022 exception.} The 2022 rate-hike episode shows reversed dynamics
(SMB$\to$HML significant). We interpret this as evidence that our finding is
specific to \emph{liquidity-driven} crises (2008, 2011, 2020) where forced
deleveraging is the dominant mechanism. Monetary policy tightening may operate
through different channels (discount-rate effects on long-duration assets),
producing different factor dynamics. We emphasize this distinction between crisis
types is speculative---we lack direct evidence for either mechanism and cannot
rule out that 2022 simply reflects sampling variation.

\textbf{Elevated regime finding.} Unlike the crisis pattern, the Elevated-regime
result (SMB$\to$HML) shows weak out-of-sample support: only 47\% of individual years
favor the expected direction, and pooled 2010--2024 data shows the opposite direction.
We therefore do not emphasize this finding and focus our conclusions on the crisis result.

\subsection{Trading Strategy Evaluation}
\label{sec:trading}

We evaluate whether the predictive relationship translates to profits via a simple
strategy: during Crisis regime, go long SMB when HML return over the past 9 days
is positive, short SMB when negative.

\begin{table}[t]
\centering
\caption{Trading Strategy Backtest (1995--2024). Strategy trades SMB based on
9-day lagged HML signal during Crisis regimes only.}
\label{tab:trading}
\begin{tabular}{lcc}
\toprule
Metric & Strategy & Buy-and-Hold SMB \\
\midrule
Annual Return & $-0.3\%$ & $+0.1\%$ \\
Sharpe Ratio & $-0.07$ & $+0.06$ \\
Max Drawdown & $-30.0\%$ & $-43.5\%$ \\
\bottomrule
\end{tabular}
\end{table}

\textbf{Why does Granger causality not imply trading profits?}

The negative returns highlight an important distinction: \emph{statistical predictability}
$\neq$ \emph{economic predictability}. Several factors explain the gap:

\begin{enumerate}
    \item \textbf{Effect size:} The Granger test detects whether HML \emph{improves}
    SMB prediction, not whether the improvement is large. The $R^2$ increment from
    adding HML lags is approximately 2\%---statistically significant but economically small.

    \item \textbf{Sign vs.\ magnitude:} Granger causality tests whether lagged HML
    coefficients are jointly non-zero, not whether they reliably predict SMB \emph{direction}.
    The relationship may involve complex dynamics (e.g., mean reversion) that a simple
    directional strategy cannot capture.

    \item \textbf{Regime detection lag:} The HMM assigns regimes using filtered
    (smoothed) probabilities. Real-time regime detection would be noisier, introducing
    additional error.

    \item \textbf{Transaction costs:} Not modeled, but would further erode returns.
\end{enumerate}

\textbf{Hedging vs.\ alpha.} The practical value, if any, lies in \emph{risk monitoring}
rather than directional trading. When the HMM signals Crisis regime and HML shows
large negative returns, a risk manager might:
\begin{itemize}
    \item Increase attention to SMB exposure
    \item Tighten stop-losses on small-cap positions
    \item Pre-arrange liquidity for potential SMB hedging
\end{itemize}

These are qualitative adjustments informed by short-horizon lead-lag evidence, not mechanical
trading rules. We do not claim demonstrated hedging value---that would require
a formal risk management backtest beyond our scope.

\subsection{Robustness}

\textbf{Lag specification.} BIC selects lag 1 in the primary regime-specific
Granger table. Normal remains strongly significant at lag 1, while Crisis is
non-significant.

\textbf{Subsample stability.} Splitting at 2008: Crisis HML$\to$SMB is
non-significant both pre-2008 ($n = 659$, $p = 0.335$) and post-2008
($n = 322$, $p = 0.599$).

\textbf{Alternative regime methods.} We tested threshold-based volatility regimes
(realized volatility $>$ 90th percentile = Crisis). These checks are exploratory
and do not overturn the primary-fit conclusion.

\textbf{Weekly aggregation.} Using weekly returns, the HML$\to$SMB crisis pattern
remains non-significant ($p = 0.114$, $n = 221$).

\textbf{Regime transitions.} Analyzing 60-day windows around Normal$\to$Crisis
transitions (determined by first day of 5+ consecutive Crisis assignments),
the HML$\to$SMB relationship emerges upon crisis entry ($p = 0.31 \to 0.04$)
and dissipates upon exit ($p = 0.02 \to 0.28$), suggesting discrete changes.

\subsection{Limitations}
\label{sec:limitations}

\begin{enumerate}
    \item \textbf{Granger $\neq$ structural causality.} Our core finding is that
    HML \emph{Granger-causes} SMB in the Normal regime---i.e., lagged HML improves SMB
    prediction conditional on lagged SMB. This does not establish that Value
    movements \emph{cause} Size movements in the interventional sense. A common
    factor (liquidity, sentiment) could drive both with different lags. The
    ``deleveraging cascade'' interpretation is a plausible hypothesis, not a verified
    mechanism.

    \item \textbf{No direct mechanism evidence.} We do not analyze holdings overlap
    (13F data), fund flows, or other direct evidence for the proposed economic channel.

    \item \textbf{Sample overlap in regime identification.} Regime labels use full-sample
    information (Section~\ref{sec:overlap}). While distributional and predictive features
    are distinct, this is not a clean train/test split.

    \item \textbf{Regime boundary exclusion.} Dropping 8\% of observations at regime
    boundaries may introduce selection bias and misses potentially interesting
    transition dynamics.

    \item \textbf{Limited crisis sample.} Crisis regime contains 1,071 days in the
    primary fit, limiting power for weaker effects.

    \item \textbf{P-value interpretation.} The reported primary $p = 8.75 \times 10^{-9}$
    reflects the specific model choices (lag selection, regime count, factor pair
    after screening). The true uncertainty is higher than this point estimate suggests.

    \item \textbf{Elevated and Crisis findings are weaker.} We avoid confirmatory
    claims outside Normal under the declared threshold.

    \item \textbf{No demonstrated hedging value.} While we suggest risk management
    applications, we do not provide a formal hedging backtest demonstrating value-add.
    The trading strategy failure indicates the predictive signal is weak.

    \item \textbf{Single factor pair.} Only HML--SMB survives multiple testing correction.
    We cannot claim broad regime-dependent structure across all factor pairs.
\end{enumerate}

\section{Conclusion}

We document regime-heterogeneous predictive evidence between HML and SMB.
The confirmatory result is a Normal-regime lag-1 HML$\to$SMB effect; Elevated
is suggestive and Crisis is null in the leakage-safe primary fit.

We emphasize what this finding does and does not establish:

\textbf{Does establish:}
\begin{itemize}
    \item A confirmatory Normal-regime HML$\to$SMB predictive relation
    \item Mixed but nonzero out-of-sample support (2/6 significant; 4/6 directional)
    \item No evidence of tradable alpha from the tested strategy
\end{itemize}

\textbf{Does not establish:}
\begin{itemize}
    \item Structural causality (common drivers could explain the pattern)
    \item Trading profitability (the relationship is too weak for directional strategies)
    \item Verified economic mechanism (deleveraging cascade is a hypothesis)
\end{itemize}

For practitioners, the finding suggests that during detected crisis regimes,
monitoring Value factor movements can provide context for SMB risk. Whether this
translates to practical hedging value remains to be demonstrated.

\begin{acks}
The author thanks the anonymous reviewers for constructive feedback that improved
the clarity of this work.
\end{acks}

\section*{Data and Code Availability}
Fama-French factor data are publicly available from Kenneth French's data library
(\url{https://mba.tuck.dartmouth.edu/pages/faculty/ken.french/data_library.html}).
Code for reproducing the analysis is available at
\href{https://github.com/ChorokLeeDev/ai-in-finance}{\nolinkurl{github.com/ChorokLeeDev/ai-in-finance}}.

\bibliography{references}

\end{document}
